\documentclass[11pt]{article}
% ===========================================================================
%  Shared preamble for all MPM2D worksheets — input right after \documentclass.
%  (Files beginning with "_" are skipped by mpm2d-worksheet-publish.mjs.)
% ===========================================================================
\usepackage[letterpaper,top=2.2cm,bottom=1.9cm,left=1.6cm,right=1.6cm,headheight=28pt,footskip=24pt]{geometry}
\usepackage{amsmath,amssymb}
\usepackage[T1]{fontenc}
\usepackage[utf8]{inputenc}
\usepackage{lmodern}
\usepackage{xcolor}
\usepackage[most]{tcolorbox}
\usepackage{enumitem}
\usepackage{multicol}
\usepackage{fancyhdr}
\usepackage{tikz}
\usetikzlibrary{calc}
\usepackage{pgfplots}
\pgfplotsset{compat=1.18}

\definecolor{exblue}{HTML}{4A90E2}
\definecolor{exbg}{HTML}{E6F3FF}
\definecolor{qorange}{HTML}{E69138}
\definecolor{qbg}{HTML}{FFF7CC}
\definecolor{keygray}{HTML}{475569}
\definecolor{keybg}{HTML}{F1F5F9}
\definecolor{iagreen}{HTML}{14653B}
\definecolor{titleblue}{HTML}{1D4ED8}
% Rotating palette so each worked example gets its own colour.
\definecolor{ex0}{HTML}{2563A0}\definecolor{ex1}{HTML}{3B7D3B}\definecolor{ex2}{HTML}{8A5A00}
\definecolor{ex3}{HTML}{A3327A}\definecolor{ex4}{HTML}{6D28A3}\definecolor{ex5}{HTML}{B08900}
\definecolor{ex6}{HTML}{0E7490}\definecolor{ex7}{HTML}{9A3412}\definecolor{ex8}{HTML}{15803D}

\pagestyle{fancy}
\fancyhf{}
\fancyhead[L]{\footnotesize NAME: \underline{\hspace{3.4cm}}}
\fancyhead[C]{\textbf{Integration Academy}\\[-2pt]{\scriptsize Math Simplified \textperiodcentered{} Success Amplified}}
\fancyhead[R]{\footnotesize MPM2D}
\fancyfoot[L]{\scriptsize dr.merie@IntegrationAcademy.ca}
\fancyfoot[C]{\scriptsize\thepage}
\fancyfoot[R]{\scriptsize IntegrationAcademy.ca}
\renewcommand{\headrulewidth}{1.2pt}
\renewcommand{\footrulewidth}{0.4pt}
\setlength{\parindent}{0pt}
\setlength{\parskip}{4pt}

% Examples are NOT breakable, so each example + its solution stays on one page.
\newtcolorbox{example}[2]{colframe=#1,colback=#1!7,coltitle=white,colbacktitle=#1,fonttitle=\bfseries,title={#2},arc=3pt,boxrule=1pt}
\newtcolorbox{qbox}[1]{colframe=qorange,colback=qbg,coltitle=white,colbacktitle=qorange,fonttitle=\bfseries,title={#1},arc=3pt,breakable,boxrule=1pt}
\newtcolorbox{keybox}{colframe=keygray,colback=keybg,coltitle=white,colbacktitle=keygray,fonttitle=\bfseries,title={$\checkmark$\ Answer Key},arc=3pt,breakable,boxrule=1pt}
\newtcolorbox{recapbox}{colframe=keygray,colback=keybg,coltitle=white,colbacktitle=keygray,fonttitle=\bfseries,title={Question Recap},arc=3pt,breakable,boxrule=1pt}
\newtcolorbox{ideabox}{colframe=titleblue,colback=titleblue!6,coltitle=white,colbacktitle=titleblue,fonttitle=\bfseries,title={The Big Ideas},arc=3pt,breakable,boxrule=1.2pt}
\newtcolorbox{learnbox}{colframe=iagreen,colback=iagreen!5,coltitle=white,colbacktitle=iagreen,fonttitle=\bfseries,title={Concept Essentials},arc=3pt,breakable,boxrule=1.2pt}
\newcommand{\lh}[1]{\par\smallskip{\bfseries\color{iagreen}#1}\par\nobreak\smallskip}

\newcommand{\soln}{\par\smallskip\textit{\textbf{Solution.}}\ }
\newcommand{\workspace}[1]{\par\vspace{3pt}\fcolorbox{black!30}{white}{\begin{minipage}[t][#1][t]{\dimexpr\linewidth-2\fboxsep\relax}\ \end{minipage}}\par}
\newcommand{\sech}[1]{\par\vspace{8pt}{\Large\bfseries\color{iagreen}#1}\par\nobreak\vspace{2pt}{\color{black!20}\hrule height1pt}\vspace{6pt}}

% A compact coordinate plot sized for an example box: \eplot{xmin}{xmax}{ymin}{ymax}{body}
\newcommand{\eplot}[5]{\begin{center}\begin{tikzpicture}
\begin{axis}[width=6.8cm,height=4.4cm,axis lines=middle,xlabel={\tiny$x$},ylabel={\tiny$y$},
grid=both,grid style={gray!16},xmin=#1,xmax=#2,ymin=#3,ymax=#4,
xtick distance=2,ytick distance=2,tick label style={font=\tiny},axis line style={-},samples=120]
#5
\end{axis}\end{tikzpicture}\end{center}}

% A sine-curve plot (degrees on x, 0..360): \splot{ymin}{ymax}{body}
\newcommand{\splot}[3]{\begin{center}\begin{tikzpicture}
\begin{axis}[width=8.6cm,height=4.6cm,axis lines=middle,xlabel={\tiny$x\,(^\circ)$},ylabel={\tiny$y$},
grid=both,grid style={gray!16},xmin=0,xmax=360,ymin=#1,ymax=#2,
xtick distance=90,ytick distance=1,tick label style={font=\tiny},axis line style={-},samples=180]
#3
\end{axis}\end{tikzpicture}\end{center}}

% A small coordinate plot: \plot{xmin}{xmax}{ymin}{ymax}{body}
\newcommand{\plot}[5]{\begin{center}\begin{tikzpicture}
\begin{axis}[width=8.4cm,height=6cm,axis lines=middle,xlabel={\scriptsize$x$},ylabel={\scriptsize$y$},
grid=both,grid style={gray!16},xmin=#1,xmax=#2,ymin=#3,ymax=#4,
xtick distance=2,ytick distance=2,tick label style={font=\tiny},axis line style={-}]
#5
\end{axis}\end{tikzpicture}\end{center}}

% Right triangle (angle theta at bottom-left, right angle at bottom-right):
%   \rtri{drawBase}{drawHeight}{oppLabel}{adjLabel}{hypLabel}
\newcommand{\rtri}[5]{\begin{center}\begin{tikzpicture}[scale=0.85]
\coordinate (A) at (0,0);\coordinate (B) at (#1,0);\coordinate (C) at (#1,#2);
\draw[thick,exblue] (A)--(B)--(C)--cycle;
\draw[black] ($(B)+(-0.32,0)$)--($(B)+(-0.32,0.32)$)--($(B)+(0,0.32)$);
\node[below] at ($(A)!0.5!(B)$){#4};
\node[right] at ($(B)!0.5!(C)$){#3};
\node[above left] at ($(A)!0.5!(C)$){#5};
\node at (0.7,0.22){$\theta$};
\end{tikzpicture}\end{center}}

% General (oblique) triangle with vertex labels and opposite-side labels:
%   \gtri{Alabel}{Blabel}{Clabel}{aSide}{bSide}{cSide}
\newcommand{\gtri}[6]{\begin{center}\begin{tikzpicture}[scale=0.85]
\coordinate (A) at (0,0);\coordinate (B) at (5,0);\coordinate (C) at (1.6,3);
\draw[thick,exblue] (A)--(B)--(C)--cycle;
\node[below left] at (A){#1};\node[below right] at (B){#2};\node[above] at (C){#3};
\node[below] at ($(A)!0.5!(B)$){#6};
\node[right] at ($(B)!0.5!(C)$){#4};
\node[left] at ($(A)!0.5!(C)$){#5};
\end{tikzpicture}\end{center}}

\fancyhead[R]{\footnotesize MCR3U \textperiodcentered{} 2.2}
\begin{document}

\begin{center}
{\LARGE\bfseries\color{titleblue}Worksheet 2.2 --- Factoring Polynomials}\\[2pt]
{\small Grade 11 Functions, University (MCR3U) \textperiodcentered{} Unit 2: Equivalent Algebraic Expressions}
\end{center}

\begin{ideabox}
Factoring is expansion in reverse: rewrite a sum as a product. Run one decision order every time, and always check by expanding.
\begin{itemize}[leftmargin=*,itemsep=2pt]
  \item \textbf{Order:} common factor first $\rightarrow$ count the terms $\rightarrow$ factor again if a bracket still factors $\rightarrow$ check by expanding.
  \item $x^2+bx+c$: find two numbers that \textbf{multiply to $c$} and \textbf{add to $b$}.
  \item $ax^2+bx+c$ with $a\neq1$: find two numbers that \textbf{multiply to $ac$} and \textbf{add to $b$}, split the middle term, then group.
  \item Patterns: $a^2-b^2=(a+b)(a-b)$ and $a^2\pm2ab+b^2=(a\pm b)^2$. A sum of squares does not factor.
  \item If no pair of factors of $ac$ adds to $b$, the trinomial is \textbf{prime}.
\end{itemize}
\end{ideabox}

\begin{learnbox}
\lh{One decision order}Run the same checklist every time: (1) take out the \textbf{greatest common factor}; (2) count the terms --- two squares subtracted is a \textbf{difference of squares}, three terms is a \textbf{trinomial}; (3) check whether any bracket factors again; (4) \textbf{check by expanding}.
\lh{Trinomials with $a=1$}$x^2+bx+c=(x+p)(x+q)$ where $p,q$ \textbf{multiply to $c$} and \textbf{add to $b$}. If $c>0$ the signs match (both take the sign of $b$); if $c<0$ the signs differ and the larger number takes the sign of $b$.
\lh{Trinomials with $a\ne1$}Expanding $(px+q)(rx+s)$ shows the two pieces of the middle term multiply to $ac$ and add to $b$. So: find $ac$; find the pair that multiplies to $ac$ and adds to $b$ (list the factor pairs in a table); \textbf{split} the middle term; \textbf{group} and pull out the common bracket.
\lh{Finishing and checking}If the brackets do not match, recheck the signs --- when a group starts with a minus, factor the minus out. If no pair of factors of $ac$ adds to $b$, the trinomial is \textbf{prime}. Always carry the common factor to the end and verify by expanding.
\end{learnbox}

\sech{Worked Examples}

\begin{example}{ex0}{Example 1 --- Common factor, then check the bracket}
Factor fully $18x^3-24x^2+6x$.\soln \textbf{Step 1 --- the GCF.} $\gcd(18,24,6)=6$ and the lowest power of $x$ is $x^1$, so the GCF is $6x$. \\ \textbf{Step 2 --- divide each term:} $\tfrac{18x^3}{6x}=3x^2,\ \tfrac{-24x^2}{6x}=-4x,\ \tfrac{6x}{6x}=1$, so $6x(3x^2-4x+1)$. \\ \textbf{Step 3 --- look inside the bracket.} $3x^2-4x+1$ has $ac=3$; the pair $-1,-3$ multiplies to $3$ and adds to $-4$: $3x^2-3x-x+1=3x(x-1)-1(x-1)=(3x-1)(x-1)$. \\ \textbf{Check:} $6x(3x-1)(x-1)=6x(3x^2-4x+1)=18x^3-24x^2+6x$. \\ \textbf{Answer:} $6x(3x-1)(x-1)$.
\end{example}

\begin{example}{ex1}{Example 2 --- Trinomial with $a=1$}
Factor $x^2+3x-28$.\soln \textbf{Need:} product $-28$, sum $+3$. One number is positive, one negative, and the larger one is positive.
\begin{center}\small\begin{tabular}{|l|c|c|c|}\hline Pair (product $-28$) & $-1,28$ & $-2,14$ & $-4,7$ \\ \hline Sum & $27$ & $12$ & $\mathbf{3}$ \\ \hline\end{tabular}\end{center}
The pair is $-4$ and $7$: $(x+7)(x-4)$. \\ \textbf{Check:} $x^2-4x+7x-28=x^2+3x-28$.
\end{example}

\begin{example}{ex2}{Example 3 --- $a\neq1$: the five steps}
Factor $3x^2+10x+8$.\soln \textbf{1. Clear the way:} no common factor. \\ \textbf{2. Find $ac$:} $3\times8=24$. \\ \textbf{3. Find the pair:} product $24$, sum $10$, all positive: $1,24\ (25)$; $2,12\ (14)$; $3,8\ (11)$; $\mathbf{4,6\ (10)}$. \\ \textbf{4. Split:} $3x^2+4x+6x+8$. \\ \textbf{5. Group:} $x(3x+4)+2(3x+4)$. Both groups give the same bracket, which confirms the split. $=(x+2)(3x+4)$. \\ \textbf{Check:} $(x+2)(3x+4)=3x^2+4x+6x+8=3x^2+10x+8$.
\end{example}

\begin{example}{ex3}{Example 4 --- $a\neq1$: a negative middle term}
Factor $5x^2-13x+6$.\soln \textbf{Find $ac$:} $30$. Product $+30$ and sum $-13$ means both numbers are \emph{negative}: $-1,-30\ (-31)$; $-2,-15\ (-17)$; $\mathbf{-3,-10\ (-13)}$. \\ \textbf{Split:} $5x^2-10x-3x+6$. \\ \textbf{Group:} $5x(x-2)-3(x-2)$. \emph{Sign trap:} the second group $-3x+6$ is written $-3(x-2)$, factoring out the $-3$ so the bracket matches. $=(5x-3)(x-2)$. \\ \textbf{Check:} $5x^2-10x-3x+6=5x^2-13x+6$.
\end{example}

\begin{example}{ex4}{Example 5 --- $a\neq1$: a negative constant}
Factor $4x^2+4x-15$.\soln \textbf{Find $ac$:} $4\times(-15)=-60$. Need product $-60$, sum $+4$: the larger number is positive, so look for two factors of $60$ whose \emph{difference} is $4$.
\begin{center}\small\begin{tabular}{|l|c|c|c|c|c|c|}\hline Factors of $60$ & $1,60$ & $2,30$ & $3,20$ & $4,15$ & $5,12$ & $6,10$ \\ \hline Difference & $59$ & $28$ & $17$ & $11$ & $7$ & $\mathbf{4}$ \\ \hline\end{tabular}\end{center}
The numbers are $+10$ and $-6$. \textbf{Split:} $4x^2+10x-6x-15$. \textbf{Group:} $2x(2x+5)-3(2x+5)=(2x-3)(2x+5)$. \\ \textbf{Check:} $4x^2+10x-6x-15=4x^2+4x-15$.
\end{example}

\begin{example}{ex5}{Example 6 --- Common factor first, then $a\neq1$}
Factor fully $18x^2+3x-6$.\soln \textbf{GCF:} all coefficients are multiples of $3$: $3(6x^2+x-2)$. \emph{This keeps $ac$ small:} $-12$ instead of $-108$. \\ \textbf{Inside:} $ac=6\times(-2)=-12$; product $-12$, sum $+1$: $1,12\ (11)$; $2,6\ (4)$; $\mathbf{3,4\ (1)}$, so $+4$ and $-3$. \\ \textbf{Split and group:} $6x^2+4x-3x-2=2x(3x+2)-1(3x+2)=(2x-1)(3x+2)$. \\ \textbf{Put the GCF back:} $3(2x-1)(3x+2)$. \\ \textbf{Check:} $(2x-1)(3x+2)=6x^2+x-2$, and $3(6x^2+x-2)=18x^2+3x-6$.
\end{example}

\begin{example}{ex6}{Example 7 --- A prime trinomial}
Can $2x^2-3x+5$ be factored over the integers?\soln $ac=10$. Product $+10$, sum $-3$, so both numbers would be negative. \\ All pairs: $-1,-10\ (-11)$; $-2,-5\ (-7)$. \\ Neither sum is $-3$, so the trinomial is \textbf{prime} --- it cannot be factored over the integers. Listing \emph{every} pair is what proves it.
\end{example}

\begin{example}{ex7}{Example 8 --- Difference of squares and perfect square}
Factor \textbf{(a)} $49x^2-36$ and \textbf{(b)} $16x^2+24x+9$.\soln \textbf{(a)} Two terms, subtracted, both squares: $(7x)^2-6^2=(7x+6)(7x-6)$. \\ \textbf{(b)} $16x^2=(4x)^2$, $9=3^2$, and $2(4x)(3)=24x$ matches the middle term, so it is a perfect square: $(4x+3)^2$. \\ \textbf{Check (b):} $(4x+3)^2=16x^2+24x+9$. (Decomposition agrees: $ac=144$, pair $12,12$.)
\end{example}

\begin{example}{ex8}{Example 9 --- A negative leading coefficient}
Factor $-6x^2+x+15$.\soln \textbf{Step 1 --- make $a$ positive:} factor out $-1$: $-(6x^2-x-15)$. \\ \textbf{Step 2:} $ac=6\times(-15)=-90$; product $-90$, sum $-1$ (larger number negative): $1,90\ (89)$; $2,45\ (43)$; $3,30\ (27)$; $5,18\ (13)$; $6,15\ (9)$; $\mathbf{9,10\ (1)}$, so $-10$ and $+9$. \\ \textbf{Step 3:} $6x^2-10x+9x-15=2x(3x-5)+3(3x-5)=(2x+3)(3x-5)$. \\ \textbf{Step 4 --- restore the $-1$:} $-(2x+3)(3x-5)$. \\ \textbf{Check:} $(2x+3)(3x-5)=6x^2-x-15$, and its negative is $-6x^2+x+15$.
\end{example}

\clearpage
\sech{Your Turn --- Show Your Work}
For each question, show all steps clearly in the space provided.

\begin{qbox}{Question 1}
Factor $15x^2-25x$.\workspace{2.2cm}
\end{qbox}

\begin{qbox}{Question 2}
Factor fully $14x^3+21x^2$.\workspace{2.2cm}
\end{qbox}

\begin{qbox}{Question 3}
Factor $x^2+11x+24$.\workspace{2.6cm}
\end{qbox}

\begin{qbox}{Question 4}
Factor $x^2-2x-35$. List the factor pairs of $35$ in a table.\workspace{2.6cm}
\end{qbox}

\begin{qbox}{Question 5}
Factor $2x^2+9x+10$. Show $ac$, the pair, the split and the grouping, then check.\workspace{3.3cm}
\end{qbox}

\begin{qbox}{Question 6}
Factor $7x^2-16x+4$. (Both numbers will be negative --- why?)\workspace{3.3cm}
\end{qbox}

\begin{qbox}{Question 7}
Factor $3x^2+x-10$. List the factor pairs of $ac$ to find the right pair.\workspace{3.3cm}
\end{qbox}

\begin{qbox}{Question 8}
Factor $6x^2-17x+5$. Watch the sign when you group.\workspace{3.3cm}
\end{qbox}

\begin{qbox}{Question 9}
Factor fully $20x^2+30x-20$.\workspace{3.6cm}
\end{qbox}

\begin{qbox}{Question 10}
Decide whether $4x^2+3x+2$ can be factored over the integers. Justify your answer.\workspace{2.6cm}
\end{qbox}

\begin{qbox}{Question 11}
Factor $64x^2-1$.\workspace{2.2cm}
\end{qbox}

\begin{qbox}{Question 12}
Factor $25x^2+20x+4$ and explain how you recognised the pattern.\workspace{2.6cm}
\end{qbox}

\begin{qbox}{Question 13 --- Challenge}
Factor fully $24x^2-10x-4$, and check your answer by expanding.\workspace{4.2cm}
\end{qbox}

\clearpage
\sech{Question Recap \& Answer Key}

\begin{recapbox}
\begin{multicols}{2}
\small
\begin{enumerate}[leftmargin=*,itemsep=3pt]
  \item Factor $15x^2-25x$.
  \item Factor fully $14x^3+21x^2$.
  \item Factor $x^2+11x+24$.
  \item Factor $x^2-2x-35$. List the factor pairs of $35$ in a table.
  \item Factor $2x^2+9x+10$. Show $ac$, the pair, the split and the grouping, then check.
  \item Factor $7x^2-16x+4$. (Both numbers will be negative --- why?)
  \item Factor $3x^2+x-10$. List the factor pairs of $ac$ to find the right pair.
  \item Factor $6x^2-17x+5$. Watch the sign when you group.
  \item Factor fully $20x^2+30x-20$.
  \item Decide whether $4x^2+3x+2$ can be factored over the integers. Justify your answer.
  \item Factor $64x^2-1$.
  \item Factor $25x^2+20x+4$ and explain how you recognised the pattern.
  \item Factor fully $24x^2-10x-4$, and check your answer by expanding.
\end{enumerate}
\end{multicols}
\end{recapbox}

\begin{keybox}
\begin{multicols}{2}
\small
\begin{enumerate}[leftmargin=*,itemsep=3pt]
  \item $5x(3x-5)$
  \item $7x^2(2x+3)$
  \item $(x+3)(x+8)$
  \item $(x-7)(x+5)$; product $-35$, sum $-2$: pairs $1,-35$ and $5,-7$ (sum $-2$)
  \item $(2x+5)(x+2)$; $ac=20$, pair $4,5$: $2x^2+4x+5x+10=2x(x+2)+5(x+2)$
  \item $(7x-2)(x-2)$; $ac=28$, pair $-14,-2$ (product positive, sum negative): $7x^2-14x-2x+4=7x(x-2)-2(x-2)$
  \item $(3x-5)(x+2)$; $ac=-30$, pair $6,-5$: $3x^2+6x-5x-10=3x(x+2)-5(x+2)$
  \item $(3x-1)(2x-5)$; $ac=30$, pair $-15,-2$: $6x^2-15x-2x+5=3x(2x-5)-1(2x-5)$
  \item $10(2x-1)(x+2)$; GCF $10$, then $2x^2+3x-2$: $ac=-4$, pair $4,-1$: $2x^2+4x-x-2=2x(x+2)-1(x+2)$
  \item prime: $ac=8$, pairs $1,8\ (9)$ and $2,4\ (6)$; neither sums to $3$
  \item $(8x-1)(8x+1)$
  \item $(5x+2)^2$; $(5x)^2$, $2^2$, and $2(5x)(2)=20x$ matches the middle term
  \item $2(4x+1)(3x-2)$; GCF $2$, then $12x^2-5x-2$: $ac=-24$, pair $-8,3$: $12x^2-8x+3x-2=4x(3x-2)+1(3x-2)$; check $(4x+1)(3x-2)=12x^2-5x-2$
\end{enumerate}
\end{multicols}
\end{keybox}

\end{document}
